\section{Arquitecturas para clasificación de imágenes}
\label{sec:marco_vision}

La red neuronal de la Sección~\ref{sec:marco_fundamentos} recibe su entrada como un vector y procesa cada valor
sin tener en cuenta dónde está ubicado. En una imagen, sin embargo, la ubicación es importante: los píxeles vecinos
suelen pertenecer a un mismo borde o a una misma región, y un mismo patrón, como el borde de una lesión, puede
aparecer en cualquier lugar. Las arquitecturas para clasificación de imágenes incorporan esta estructura de dos
maneras, la convolución y la atención. Los cuatro clasificadores que evalúa el trabajo pertenecen a esas dos
familias: DenseNet y ConvNeXt son redes convolucionales, y el \gls{vit} y el Swin Transformer se basan en la
atención. La sección presenta estas arquitecturas comenzando con las redes convolucionales y siguiendo con la
atención, el transformer y sus adaptaciones a imágenes.

\subsection{Redes convolucionales}
\label{sec:marco_cnn}

Las \glspl{cnn} parten de la idea de mirar la imagen por partes pequeñas en lugar de toda junta. La operación
que lo hace es la \textbf{convolución} (\emph{convolution}), y su herramienta es un filtro (\emph{kernel}): una
pequeña matriz de pesos, por ejemplo de $3 \times 3$, que codifica un patrón visual, como el paso de un tono
claro a uno oscuro. El filtro se apoya sobre una ventana de $3 \times 3$ píxeles, multiplica cada peso por el
píxel que tiene debajo y suma los nueve productos, es decir, calcula el producto punto entre sus pesos y la
ventana. Cuanto más se parece la ventana al patrón que codifica el filtro, mayor es el resultado. Por ejemplo,
un filtro con pesos positivos en su columna izquierda, nulos en la central y negativos en la derecha da un valor
grande sobre una ventana clara a la izquierda y oscura a la derecha, porque los píxeles claros, que tienen
valores mayores que los oscuros, se multiplican por los pesos positivos. Sobre una ventana de tono uniforme, en
cambio, las dos columnas se cancelan y el resultado es cero. Después, el filtro se desplaza un píxel y repite el cálculo, hasta recorrer toda la imagen. Los resultados, ordenados según la
posición en que se calcularon, forman un mapa de características (\emph{feature map}) que indica dónde aparece
el patrón. En una imagen dermatoscópica, ese filtro produce valores grandes en los tramos verticales del contorno
de la lesión, donde la piel pasa de clara a oscura, y valores cercanos a cero sobre la piel sana, de tono
uniforme.

% Figura: la operación de convolución.
% Se usa en sec:marco_cnn, después del párrafo que define la convolución.
% Las tres matrices van separadas y las etiquetas en dos líneas, para que no se toquen (tutor, p36).

\begin{figure}[H]
    \centering
    \begin{tikzpicture}[x=1mm, y=1mm]

        % Entrada 6x6 con la ventana de 3x3 resaltada.
        \fill[fignaranja!30] (5, 10) rectangle (20, 25);
        \draw[figtrazo!50] (0, 0) grid[step=5] (30, 30);
        \draw[figtrazo, thick] (0, 0) rectangle (30, 30);
        \draw[fignaranja!70!black, very thick] (5, 10) rectangle (20, 25);
        \node[font=\footnotesize, text=figtrazo, align=center, anchor=north] at (15, -2) {entrada\\$6\times 6$};

        % Filtro 3x3.
        \fill[fignaranja!30] (50, 7.5) rectangle (65, 22.5);
        \draw[fignaranja!70!black, shift={(50, 7.5)}] (0, 0) grid[step=5] (15, 15);
        \draw[fignaranja!70!black, very thick] (50, 7.5) rectangle (65, 22.5);
        \node[font=\footnotesize, text=figtrazo, align=center, anchor=north] at (57.5, -2) {filtro $3\times 3$\\(pesos compartidos)};
        \node[font=\large, text=figtrazo] at (40, 15) {$\ast$};

        % Salida 4x4 con la celda resultante.
        \fill[figverde!35] (95, 10) rectangle (100, 15);
        \draw[figtrazo!50] (90, 0) grid[step=5] (110, 20);
        \draw[figtrazo, thick] (90, 0) rectangle (110, 20);
        \draw[figverde!70!black, very thick] (95, 10) rectangle (100, 15);
        \node[font=\footnotesize, text=figtrazo, align=center, anchor=north] at (100, -2) {mapa de características\\$4\times 4$};
        \node[font=\large, text=figtrazo] at (77.5, 15) {$=$};

        % De la ventana a la celda.
        \draw[figflecha, fignaranja!70!black, bend left=25] (20, 25) to
            node[figetiqueta, above] {producto punto} (95, 15);

    \end{tikzpicture}
    \caption{Convolución: un mismo filtro recorre la entrada y produce un mapa de características.}
    \label{fig:marco_convolucion}
\end{figure}


La Figura~\ref{fig:marco_convolucion} muestra el recorrido sobre una entrada de $6 \times 6$: como el filtro
cabe en cuatro posiciones por fila y cuatro por columna, el mapa es de $4 \times 4$. La operación responde a los dos rasgos de las imágenes señalados antes. Como cada valor del mapa depende solo de
una ventana pequeña, la red supone que la información relevante es local, y como el mismo filtro, con los mismos
pesos, se aplica en toda la imagen, el contorno de la lesión se detecta en cualquier lugar en que aparezca. Los
pesos de los filtros no se fijan a mano, sino que se aprenden durante el entrenamiento, como los de cualquier
otra capa. Una capa convolucional aplica muchos filtros a la vez, cada uno con sus propios pesos, y apila sus
mapas como canales de la salida~\cite{lecun1998gradient,goodfellow2016deep}.

Una \gls{cnn} completa alterna capas convolucionales, funciones de activación y capas de reducción
(\emph{pooling}), que resumen cada ventana de un mapa en un solo valor, por ejemplo su máximo, y reducen así la
resolución. Cada reducción hace que una unidad de la capa siguiente dependa de una región más grande de la
imagen original, llamada su campo receptivo (\emph{receptive field}). De este modo, la red concreta la cadena
de rasgos simples a complejos descrita al comienzo de la Sección~\ref{sec:marco_fundamentos}: las primeras capas
responden a bordes y las últimas, a partes de objetos~\cite{goodfellow2016deep}. En una imagen dermatoscópica,
esa cadena iría del contorno y los cambios de color a estructuras que ocupan parte de la lesión, como una red de
líneas pigmentadas o un grupo de puntos~\cite{argenziano1998checklist}, y de ellas a la lesión completa. Al
final de la red, el
promedio de cada canal sobre todas las posiciones convierte los mapas en un vector de características, y la
cabeza de clasificación de la Sección~\ref{sec:marco_transferencia} lo convierte en los
logits~\cite{he2015residual}.

Agregar capas, sin embargo, no siempre mejora la red. A partir de cierta profundidad aparece la degradación
(\emph{degradation}): el error aumenta al sumar capas. Por ejemplo, una red convolucional de 56 capas obtiene
más error de entrenamiento que una de 20 sobre el mismo conjunto~\cite{he2015residual}. Como el error crece
sobre los propios datos de entrenamiento, el problema no es el sobreajuste de la
Sección~\ref{sec:marco_regularizacion}, sino la dificultad de optimizar la red. En principio, a la red de 56
capas le bastaría con que sus 36 capas adicionales calcularan la identidad, es decir, devolvieran su entrada
sin cambios, para igualar a la de 20. Sin embargo, al optimizador le resulta difícil aproximar la identidad con
varias capas no lineales~\cite{he2015residual}.

ResNet resolvió esta dificultad con las \textbf{conexiones residuales} (\emph{residual connections}). Una
conexión residual suma la entrada de un bloque de capas al resultado de esas capas, de modo que la entrada
llega a la salida del bloque también por un camino que no la transforma. Así, el bloque ya no tiene que producir su salida completa, sino solo la
corrección $F$ que hay que agregarle a la entrada,
\begin{equation}
    \mathbf{h}^{(l+1)} = \mathbf{h}^{(l)} + F\bigl(\mathbf{h}^{(l)}\bigr),
    \label{eq:marco_residual}
\end{equation}
donde $\mathbf{h}^{(l)}$ es la entrada del bloque $l$. Si lo mejor para un bloque es no modificar su entrada, le
basta con llevar los pesos de $F$ a cero, y su salida es igual a su entrada. De este modo, en el peor caso un
bloque nuevo no cambia nada, y en el mejor aporta una corrección útil. Con esta modificación se pudieron entrenar redes de más de cien capas~\cite{he2015residual}, y la misma conexión
reaparece en el transformer (Sección~\ref{sec:marco_atencion}).

Las dos redes convolucionales que evalúa el trabajo parten de esta idea y la desarrollan en direcciones
distintas. DenseNet conecta cada capa con todas las anteriores de su bloque: cada capa recibe, concatenados, los
mapas de características de las capas previas, de modo que reutiliza los rasgos ya calculados y alcanza una
exactitud comparable a la de una ResNet con menos parámetros~\cite{huang2018densely}. ConvNeXt, en cambio,
conserva la estructura de una ResNet y actualiza su diseño con decisiones tomadas de los transformers que se
presentan a continuación, como filtros de $7 \times 7$ y la normalización por capa~\cite{liu2022convnet}.
Aun con estas mejoras, una convolución solo relaciona píxeles cercanos. Para vincular dos regiones alejadas de
la imagen, por ejemplo, para comparar un borde de la lesión con el opuesto, la red necesita apilar capas hasta
que su campo receptivo cubra a ambos.

\subsection{Atención y arquitectura transformer}
\label{sec:marco_atencion}

La \textbf{atención} (\emph{attention}), en cambio, relaciona cualquier par de posiciones de la entrada en una
sola capa, sin importar la distancia entre ellas. Se propuso en 2017 para la traducción de textos, como núcleo
de una arquitectura llamada transformer~\cite{vaswani2017attention}, y hoy los modelos de lenguaje y buena
parte de los modelos de imagen se construyen sobre ella~\cite{dosovitskiy2021image}. Para aplicarla, la entrada
se divide en una secuencia de elementos llamados \emph{tokens}. En un texto, cada token es una palabra o
un fragmento de palabra, y en una imagen, un parche (\emph{patch}), es decir, un pequeño recuadro de píxeles.
Cada token se representa con un \emph{embedding}, un vector numérico que el modelo aprende de modo que
los tokens con significados parecidos tengan vectores parecidos.

La atención determina, para cada token, qué información tomar de los demás. Para ello, proyecta el
embedding de cada token en tres vectores, con matrices de pesos aprendidas: una consulta
(\emph{query}), que expresa qué información busca el token; una clave (\emph{key}), que expresa qué
información ofrece a los demás, y un valor (\emph{value}), que es la información que entrega. La compatibilidad
entre dos tokens se mide con el producto punto entre la consulta de uno y la clave del otro, y la
softmax de la Ecuación~\ref{eq:marco_softmax} convierte esas compatibilidades en pesos que suman uno. La salida
de cada token es el promedio de los valores de todos los tokens, ponderado por esos pesos. Si se
reúnen las consultas, las claves y los valores en las matrices $Q$, $K$ y $V$, la operación completa, llamada
atención escalada por producto punto, es
\begin{equation}
    \mathrm{Attention}(Q, K, V) = \mathrm{softmax}\!\left(\frac{QK^{\top}}{\sqrt{d_k}}\right) V,
    \label{eq:marco_atencion}
\end{equation}
donde $d_k$ es la dimensión de las claves. La división por $\sqrt{d_k}$ evita que los productos punto crezcan
con la dimensión y lleven la softmax a zonas donde su gradiente es casi nulo, lo que dificultaría el
entrenamiento~\cite{vaswani2017attention}. La Figura~\ref{fig:marco_atencion} muestra los pesos para una
secuencia de cuatro tokens: cada fila corresponde a una consulta, y sus pesos indican cuánto aporta el
valor de cada posición a la salida de ese token.

% Figura: la matriz de pesos de atención.
% Se usa en sec:marco_atencion, después de la Ecuación de la atención escalada por producto punto.
% Los pesos son ilustrativos: muestran la forma de la operación, no un modelo entrenado.

\begin{figure}[H]
    \centering
    \begin{tikzpicture}[x=1mm, y=1mm]

        % Consultas (filas) y claves (columnas).
        \foreach \i in {1, 2, 3, 4} {
            \node[fignodo, minimum height=6mm, minimum width=9mm, inner sep=1pt,
                  draw=figazul!80!black, fill=figazul!12] (q\i) at (0, 30 - \i*8) {\footnotesize $q_{\i}$};
            \node[fignodo, minimum height=6mm, minimum width=6.5mm, inner sep=1pt,
                  draw=fignaranja!70!black, fill=fignaranja!18] (k\i) at (8 + \i*8, 34) {\footnotesize $k_{\i}$};
        }

        % Pesos softmax(QK^T / sqrt(d_k)); cada fila suma uno.
        \foreach \row/\vals in {1/{60,20,10,10}, 2/{15,55,20,10}, 3/{10,25,50,15}, 4/{5,10,25,60}} {
            \foreach \v [count=\col] in \vals {
                \fill[figceleste!\v!white] (4 + \col*8, 26 - \row*8) rectangle (12 + \col*8, 34 - \row*8);
                \node[font=\scriptsize] at (8 + \col*8, 30 - \row*8) {0.\ifnum\v<10 0\fi\v};
            }
        }
        \draw[figtrazo, shift={(12, -6)}] (0, 0) grid[step=8] (32, 32);

        % Salida.
        \foreach \i in {1, 2, 3, 4} {
            \node[fignodo, minimum height=6mm, minimum width=9mm, inner sep=1pt,
                  draw=figverde!70!black, fill=figverde!10, densely dashed] (o\i) at (66, 30 - \i*8)
                  {\footnotesize $o_{\i}$};
            \draw[figflecha] (45, 30 - \i*8) -- (o\i.west);
        }

        \node[font=\footnotesize, text=figtrazo] at (66, 32) {$o_i = \sum_j a_{ij} v_j$};
        \node[font=\footnotesize, text=figtrazo, align=center] at (28, -12)
            {pesos $a_{ij}$ = $\mathrm{softmax}\!\left(q_i k_j^{\top} / \sqrt{d_k}\right)$};

    \end{tikzpicture}
    \caption{Pesos de atención entre las posiciones de una secuencia de cuatro tokens (valores ilustrativos).}
    \label{fig:marco_atencion}
\end{figure}


Un solo cálculo de atención reúne en un único conjunto de pesos todas las relaciones entre tokens. La
atención multicabezal (\emph{multi-head attention}) realiza varios cálculos en paralelo, cada uno con sus
propias proyecciones de menor dimensión, y concatena los resultados, de modo que cada cabeza puede capturar un
tipo distinto de relación entre posiciones~\cite{vaswani2017attention}.

La atención, tal como se definió, no tiene en cuenta el orden de los tokens: si se reordena la entrada,
cada token obtiene el mismo resultado que antes. Para que el modelo distinga las posiciones, se suma al
embedding de cada token un vector que depende de su lugar en la secuencia, llamado codificación de
posición (\emph{positional encoding}), que puede ser fijo o aprenderse durante el
entrenamiento~\cite{vaswani2017attention,dosovitskiy2021image}.

Un transformer apila bloques iguales. Cada bloque aplica primero una capa de atención multicabezal, que combina
información entre tokens, y después un perceptrón multicapa, que transforma cada token por
separado. Las dos partes llevan una conexión residual como la de la Ecuación~\ref{eq:marco_residual} y una
normalización por capa (\emph{layer normalization}), que lleva los valores de cada token a una escala
común~\cite{vaswani2017attention}. El transformer original combina dos pilas de bloques: un codificador
(\emph{encoder}), que procesa la entrada completa y produce una representación de cada token que
incorpora información de los demás, y un decodificador (\emph{decoder}), que genera la salida de a un
token por vez y se presenta en la Sección~\ref{sec:marco_llm}. El codificador no depende de qué
representen los tokens, de modo que también puede procesar una imagen si sus tokens son parches.

\subsection{Vision Transformer}
\label{sec:marco_vit}

Hasta la aparición del \gls{vit}, la atención se usaba en visión junto con las \glspl{cnn} o en reemplazo de
algunas de sus partes, y la convolución seguía siendo la base de los clasificadores. El \gls{vit} probó si un
transformer sin convoluciones, aplicado directamente a la imagen, podía clasificarla~\cite{dosovitskiy2021image}.
Para ello, aplica el codificador del transformer y convierte la imagen en una secuencia de tokens
(Figura~\ref{fig:marco_vit_patches}). Primero divide la imagen en una
rejilla de parches cuadrados, por ejemplo de $16 \times 16$ píxeles, de modo que una imagen de
$224 \times 224$ píxeles da 196 parches. En una imagen dermatoscópica, cada parche cubre un fragmento de la
lesión o de la piel que la rodea. Después aplana cada parche en un vector y lo proyecta, con una capa
lineal, a un embedding de dimensión fija, al que suma una codificación de posición aprendida que indica
su lugar en la rejilla. Al comienzo de la secuencia agrega un token de clase, que no corresponde a ningún
parche. Al atravesar los bloques del codificador, ese token reúne información de todos los parches, y la
cabeza de clasificación toma su representación final para producir los
logits~\cite{dosovitskiy2021image}.

% Figura 2.1. Del parche al token en el Vision Transformer.
% Se usa en sec:marco_vision, subsección "Vision Transformer y Swin Transformer", después del
% párrafo que define el ViT. Cada parche de la fila superior y su token comparten color, para que se
% lea la correspondencia; el token de clase va en gris porque no viene de ningún parche.

\begin{figure}[H]
    \centering
    \begin{tikzpicture}[x=1mm, y=1mm]

        % Imagen dividida en una rejilla de 4x4 parches; la fila superior va coloreada.
        \draw[figtrazo, thick, fill=figrelleno] (-26, -12) rectangle (0, 12);
        \foreach \c [count=\i from 0] in {figazul, fignaranja, figverde, figpurpura} {
            \fill[\c!35] (-26 + \i*6.5, 6) rectangle (-19.5 + \i*6.5, 12);
        }
        \foreach \x in {-19.5, -13, -6.5} { \draw[figtrazo] (\x, -12) -- (\x, 12); }
        \foreach \y in {-6, 0, 6} { \draw[figtrazo] (-26, \y) -- (0, \y); }
        \draw[figtrazo, thick] (-26, -12) rectangle (0, 12);
        \node[font=\small, text=figtrazo, align=center] at (-13, -18) {Imagen dividida\\en parches};

        % Operación de proyección.
        \node[figtarea, text width=28mm] (proj) at (22, 0)
            {Proyectar cada parche\\a un token};

        % Token de clase y tokens de los parches, con su embedding de posición.
        \node[fignodo, minimum width=7mm, minimum height=6mm, inner sep=1pt, fill=figrelleno]
            (t0) at (50, 6) {\footnotesize[CLS]};
        \node[font=\footnotesize, text=figtrazo] at (50, -3) {$+$};
        \node[font=\footnotesize, text=figtrazo] at (50, -9) {$p_{0}$};
        \foreach \c [count=\i from 1] in {figazul, fignaranja, figverde, figpurpura} {
            \node[fignodo, draw=\c!70!black, fill=\c!35, minimum width=6mm, minimum height=6mm,
                  inner sep=1pt] (t\i) at (50 + \i*9, 6) {\footnotesize $x_{\i}$};
            \node[font=\footnotesize, text=figtrazo] at (50 + \i*9, -3) {$+$};
            \node[font=\footnotesize, text=figtrazo] at (50 + \i*9, -9) {$p_{\i}$};
        }
        \node[font=\small, text=figtrazo] at (95, 6) {$\cdots$};
        \node[font=\small, text=figtrazo, align=center] at (72, -18)
            {Secuencia de tokens con\\embedding de posición por parche};

        % Flechas.
        \draw[figflecha] (0, 0) -- (proj.west);
        \draw[figflecha] (proj.east) -- (43, 0) |- (t0.west);

    \end{tikzpicture}
    \caption{Del parche al token en el Vision Transformer, con el token de clase ([CLS]) al inicio.}
    \label{fig:marco_vit_patches}
\end{figure}


Esta forma de procesar la imagen supone mucho menos sobre ella que una \gls{cnn}. Los supuestos sobre los datos
que una arquitectura incorpora de antemano, antes de ver ningún ejemplo, se llaman \textbf{sesgo inductivo}
(\emph{inductive bias}). En una \gls{cnn}, la localidad y la equivariancia a traslaciones, es decir, que un
patrón desplazado produzca la misma respuesta desplazada, están presentes en cada capa. En el \gls{vit}, en
cambio, solo el perceptrón de cada bloque es local, porque opera sobre cada token por separado, mientras
que la atención relaciona todos los parches desde la primera capa. Por lo tanto, las relaciones espaciales entre
parches tienen que aprenderse de los datos, y el \gls{vit} necesita más ejemplos que una \gls{cnn}: preentrenado
solo sobre \emph{ImageNet}, alcanza una exactitud algo menor que la de \glspl{cnn} de tamaño comparable, y
preentrenado sobre conjuntos de entre 14 y 300 millones de imágenes, las supera~\cite{dosovitskiy2021image}. En
consecuencia, el desempeño de un \gls{vit} depende en gran medida del preentrenamiento del que parte antes del
ajuste fino (Sección~\ref{sec:marco_transferencia}).

\subsection{Swin Transformer}
\label{sec:marco_swin}

El \gls{vit} mostró que la atención alcanza para clasificar imágenes, pero dejó dos problemas. El primero es el
costo: la atención global crece con el cuadrado del número de parches, porque relaciona cada token con todos
los demás, y eso limita la resolución de las imágenes que se pueden procesar, ya que duplicar el lado de una
imagen cuadruplica la cantidad de parches. El segundo es que procesa toda la imagen a una sola escala: sus
parches conservan el mismo tamaño en todos los bloques, sin la pirámide de mapas cada vez más pequeños de una
\gls{cnn}, que sirve cuando los objetos aparecen a tamaños distintos~\cite{liu2021swin}. El Swin Transformer
resuelve ambos problemas.

Para reducir el costo, el Swin Transformer limita la atención a ventanas locales de $M \times M$ parches. Si el
mapa tiene $h \times w$ parches y $D$ canales, el costo de la atención global y el de la atención por ventanas
son
\begin{equation}
    \Omega_{\text{global}} = 4hwD^2 + 2(hw)^2D, \qquad \Omega_{\text{ventanas}} = 4hwD^2 + 2M^2hwD,
    \label{eq:marco_swin_costo}
\end{equation}
de modo que, con $M$ fijo, el segundo crece de forma lineal con el número de parches~\cite{liu2021swin}.

% Figura: ventanas y ventanas desplazadas del Swin Transformer. Adaptada de liu2021swin (Figura 2).
% Se usa en sec:marco_vit, en el párrafo del Swin.

\begin{figure}[H]
    \centering
    \begin{tikzpicture}[x=1mm, y=1mm]

        % Capa l: cuatro ventanas de 4x4 parches.
        \foreach \c/\x/\y in {figazul/0/16, fignaranja/16/16, figverde/0/0, figpurpura/16/0} {
            \fill[\c!22] (\x, \y) rectangle (\x + 16, \y + 16);
        }
        \draw[figtrazo!40] (0, 0) grid[step=4] (32, 32);
        \draw[figtrazo, very thick] (0, 0) rectangle (32, 32);
        \draw[figtrazo, very thick] (16, 0) -- (16, 32) (0, 16) -- (32, 16);
        \node[font=\footnotesize, text=figtrazo] at (16, -5) {capa $l$: ventanas regulares};

        % Capa l+1: ventanas desplazadas dos parches.
        \begin{scope}[xshift=50mm]
            \foreach \c/\x/\y/\w/\h in {
                figazul/0/24/8/8, fignaranja/8/24/16/8, figverde/24/24/8/8,
                figpurpura/0/8/8/16, figceleste/8/8/16/16, figbermellon/24/8/8/16,
                figazul/0/0/8/8, fignaranja/8/0/16/8, figverde/24/0/8/8} {
                \fill[\c!22] (\x, \y) rectangle (\x + \w, \y + \h);
            }
            \draw[figtrazo!40] (0, 0) grid[step=4] (32, 32);
            \draw[figtrazo, very thick] (0, 0) rectangle (32, 32);
            \draw[figtrazo, very thick] (8, 0) -- (8, 32) (24, 0) -- (24, 32) (0, 8) -- (32, 8) (0, 24) -- (32, 24);
            \node[font=\footnotesize, text=figtrazo] at (16, -5) {capa $l+1$: ventanas desplazadas};
        \end{scope}

        \draw[figflecha] (36, 16) -- (46, 16);

    \end{tikzpicture}
    \caption{Ventanas y ventanas desplazadas del Swin Transformer. Adaptada de \textcite{liu2021swin}.}
    \label{fig:marco_swin}
\end{figure}


Las ventanas fijas, sin embargo, no intercambian información entre sí. Para que la información pase de una
ventana a otra, los bloques alternan una partición regular en ventanas con otra desplazada en medio tamaño de
ventana, cuyas ventanas abarcan parches de varias ventanas de la partición anterior
(Figura~\ref{fig:marco_swin}). Para resolver el segundo problema, entre una etapa y la siguiente, una capa de
fusión de parches (\emph{patch merging})
combina cada grupo de $2 \times 2$ parches vecinos en uno solo, de modo que la resolución se reduce a la mitad
y la cantidad de canales se duplica. El resultado es una pirámide de mapas de características como la de una
\gls{cnn}. Por último, la posición de los parches no se suma al embedding como en el \gls{vit}, sino
que entra como un término que depende de la posición relativa entre cada par de parches de la ventana y se suma
a su compatibilidad~\cite{liu2021swin}.

Las arquitecturas presentadas en esta sección se diferencian, en definitiva, por cuánto suponen sobre la
imagen. Las \glspl{cnn} incorporan la localidad y la jerarquía en su diseño; el \gls{vit} prescinde de esos
supuestos y compensa su ausencia con más datos de preentrenamiento, y el Swin Transformer recupera la localidad
y la jerarquía con una operación de atención. La diferencia entre familias, sin embargo, no reside solo en la
operación de base: ConvNeXt, una red puramente convolucional con un diseño actualizado, alcanza en
\emph{ImageNet} una exactitud igual o superior a la del Swin Transformer~\cite{liu2022convnet}. Por eso el
trabajo evalúa representantes de ambas familias.

Estos clasificadores cumplen, además, un papel de referencia. Una \textbf{línea base} (\emph{baseline}) es un
modelo establecido contra el que se compara una propuesta nueva, porque un resultado aislado no permite juzgar
el desempeño de un modelo hasta que se lo contrasta con lo que logra la referencia sobre los mismos datos. Las
\glspl{cnn} y los transformers visuales entrenados con las etiquetas de la tarea son la línea base de los
\glspl{vlm} (Sección~\ref{sec:marco_vlm}): son las arquitecturas consolidadas para clasificar imágenes, y un
modelo de otra familia justifica su costo si iguala su desempeño o si aporta algo que ellas no ofrecen. Sin
embargo, ninguno de estos clasificadores produce texto, ya que su salida es una probabilidad por clase. Para
redactar un reporte hace falta un modelo capaz de generar secuencias de palabras, y los modelos de lenguaje
actuales se construyen sobre el decodificador del transformer.
